\documentclass[11pt]{article}
\usepackage[a4paper,margin=1in]{geometry}
\usepackage{amsmath,amssymb,amsthm}
\usepackage{booktabs}
\usepackage{graphicx}
\usepackage{xcolor}
\usepackage{hyperref}
\hypersetup{colorlinks=true,linkcolor=blue,urlcolor=blue,citecolor=blue}
\newcommand{\rt}[1]{\texttt{#1}}
\newcommand{\fn}[1]{\texttt{#1}}
\newcommand{\wcc}{\mathrm{wcc}}
\newcommand{\att}{\mathrm{att}}
\newcommand{\ifgraphic}[2]{\IfFileExists{#1}{\includegraphics[width=#2]{#1}}%
  {\fbox{\parbox{#2}{\centering\small figure \texttt{\detokenize{#1}} pending}}}}
\newcommand{\iftable}[1]{\IfFileExists{#1}{\input{#1}}%
  {\fbox{\parbox{0.9\textwidth}{\centering\small table
  \texttt{\detokenize{#1}} pending}}}}

\theoremstyle{plain}
\newtheorem{proposition}{Proposition}

\title{\textbf{Report R16 --- the rules that sit under the $ab\ge2$ hyperbola,
and why}\\
\large Krylov sector analysis of quantum cellular automata}
\author{Tier 1 analysis}
\date{2026-08-07}

\begin{document}
\maketitle

\begin{abstract}
A handful of markers in the left panel of R9's Fig.~F1 lie below the exclusion
curve $ab=2$, which the partition of the computational basis forbids. This report
identifies them and finds \textbf{two unrelated causes}. The larger effect is a
\emph{rendering artefact}: F1 nudges the three families sideways by $\pm0.019$ in
$a$ so that none hides under another, and because $b=2/a$ rises as $a$ falls, the
$40$ V-free rules sitting exactly \emph{on} the curve at $(1,2)$ are drawn where
the curve is $2.0387$ --- under it. They are the biggest marker in the panel. The
smaller effect is real arithmetic: \textbf{seven} rules --- $W_6$, $W_{14}$,
$W_{20}$, $W_{74}$, $W_{84}$, $W_{88}$, $W_{229}$ --- had a fitted product
$ab=1.919$--$1.937$. The answer to ``is it simply a fitting problem?'' is
\emph{yes, but the fix is a theorem rather than a better fit}. Their sector
counts are exact quasi-linear staircases, $\lceil N/2\rceil+1$ and
$\lfloor(N+4)/3\rfloor$, verified at every $N\le22$; so $a=1$, pigeonhole gives
$D_{\max}\ge2^N/n_\wcc$ and hence $b\ge2$, while $D_{\max}\le2^N$ gives $b\le2$.
The pair $(1,1.92)$ is not a violated theorem but an \emph{internally
inconsistent descriptor}. No better estimator rescues the fit on this window: the
certified control $W_{134}$, whose $D_{\max}$ \emph{is} the central binomial and
whose base is $2$ by derivation, is itself fitted at $1.9244$ and has a rolling
base of $1.904$--$1.938$, squarely inside the band spanned by the seven. We
extended the sweep to $N=22$, installed derived $(a,b)=(1,2)$ overrides with a
refitted prefactor power, and the sector map now reports $204$ rules on the
curve, $50$ above, $2$ irregular, and \textbf{zero} below.
\end{abstract}

\tableofcontents

\section{The question}

R9~\S``The hyperbola'' derives the only inequality in this project that
constrains the sector map from outside the data. The weak components partition
the $2^N$ basis states, so the largest one cannot be smaller than the average:
\begin{equation}
n_\wcc(N)\,D_{\max}(N)\;\ge\;2^N ,
\label{eq:finite}
\end{equation}
and writing $n_\wcc\sim a^N$, $D_{\max}\sim b^N$ and taking $N$-th roots,
\begin{equation}
a\,b\;\ge\;2 .
\label{eq:hyper}
\end{equation}
Equality holds when the sectors are comparable in size. The curve is drawn in the
left panel of R9's F1 and it excludes the region below it. Some markers are down
there. This report asks which, and why.

The short answer is that ``some markers are down there'' conflates two things
that have nothing to do with each other, and separating them is most of the work.

\section{Cause one: a rendering artefact, and it is the biggest marker}
\label{sec:artefact}

F1 stacks $256$ rules into few cells --- $77\%$ of the sector map is a single
point --- so \fn{sector\_figure} offsets the three families horizontally to stop
one hiding under another:
\[
\Delta a = \begin{cases}
+0.019 & \text{V+reset}\\
\phantom{+}0.000 & \text{unitary}\\
-0.019 & \text{V-free.}
\end{cases}
\]
This is documented in F1's caption as a caveat about reading absolute bases off
the axis. At the top-left corner it is more than that. The exclusion curve is
\emph{steep} there: $\mathrm{d}b/\mathrm{d}a=-2/a^2=-2$ at $a=1$, so a leftward
nudge of $0.019$ raises the curve by $0.0387$ while leaving the marker at
$b=2$. Every V-free rule at exactly $(1,2)$ --- one sector of size $2^N$, on the
curve to machine precision --- is therefore \emph{drawn} at $(0.981,2.000)$ with
the curve passing above it at $2.0387$.

There are $40$ such rules after the correction of \S\ref{sec:seven} ($35$ before
it; five of the seven are V-free and join this cell). They are the single largest
marker in the panel, and they account for more visual ``violation'' than the real
effect does. The verdict table has always reported them correctly --- their
product is $2.000000$ --- so nothing downstream was ever wrong; only the picture
was misleading. F1's caption and figure header now say so.

\section{Cause two: seven rules whose fitted product really was below $2$}
\label{sec:seven}

\subsection{Who they are}

Lifting the overrides described below and re-running the $256$-rule fit inside
the uniform window $N\le16$ gives exactly seven rules with $ab<2$:

\begin{center}
\begin{tabular}{llll}
\toprule
group & rules & tuples & family \\
\midrule
A & $W_6$, $W_{14}$, $W_{20}$, $W_{84}$ & \rt{IVDD}, \rt{IEDD}, \rt{IDVD},
  \rt{IDED} & $6,20$ mixed; $14,84$ V-free \\
B & $W_{74}$, $W_{88}$, $W_{229}$ & \rt{DEID}, \rt{DIED}, \rt{EDIE} & all V-free \\
\bottomrule
\end{tabular}
\end{center}

\noindent The grouping is not cosmetic: within each group the rules share an
\emph{identical} $n_\wcc$ series and an identical $D_{\max}$ series at every $N$
computed, so there are four distinct series among the seven, not seven. Group A
is fitted at $b=1.922559$ and group B at $1.932218$ ($W_{74}$), $1.936760$
($W_{88}$), $1.918755$ ($W_{229}$); all seven have $a=1$.

\subsection{The sector counts are exact, not fitted}

\begin{proposition}
For every $N$ on record ($6\le N\le22$, computed at open boundaries),
\[
n_\wcc(N)=\Big\lceil \tfrac{N}{2}\Big\rceil+1 \quad\text{(group A)},
\qquad
n_\wcc(N)=\Big\lfloor \tfrac{N+4}{3}\Big\rfloor \quad\text{(group B)} .
\]
\end{proposition}

\noindent This is checked exhaustively rather than fitted, and it is also
recovered independently by a residue-class solver
(\fn{below\_curve.quasilinear\_formula}) that looks for an affine form in $N$ on
each class modulo $m\le4$ using exact rational arithmetic; it returns period $2$
with common slope $1/2$ for group A and period $3$ with common slope $1/3$ for
group B. So $n_\wcc=\Theta(N)$ and, unambiguously, $a=1$.

\subsection{The squeeze}

\begin{proposition}[$b=2$ exactly]
If $n_\wcc(N)=\Theta(N)$ then
\[
b \;=\; \lim_{N\to\infty} D_{\max}(N)^{1/N} \;=\; 2 .
\]
\end{proposition}

\begin{proof}
Pigeonhole on the partition gives $D_{\max}\ge 2^N/n_\wcc \ge 2^N/(cN)$ for some
constant $c$, hence $b\ge\lim (2^N/cN)^{1/N}=2$. A single sector cannot exceed
the whole basis, so $D_{\max}\le2^N$ and $b\le2$.
\end{proof}

\noindent Both halves are elementary; the point is that \emph{once $a=1$ is
established, $b$ is not a free parameter at all}. The fitted pair $(1,1.92)$ does
not violate \eqref{eq:hyper} so much as contradict itself: the same partition
argument that draws the curve fixes the second coordinate given the first.

It is worth stating plainly why the finite-$N$ version of the bound never bites.
Equation \eqref{eq:finite} holds comfortably --- the smallest value of
$n_\wcc D_{\max}/2^N$ across the seven is $2.19$ (group A), $1.31$ ($W_{74}$,
$W_{88}$) and $1.59$ ($W_{229}$), and across all $256$ rules the tightest case
anywhere is $W_0$ at $N=6$ with ratio exactly $1.0000$. But the induced bound on
the base, $(2^N/n_\wcc)^{1/N}$, converges to $2$ at the glacial rate
$(cN)^{-1/N}$: at $N=22$ it is still only $1.79$--$1.82$. The theorem is exact
and the finite-$N$ shadow of it is useless. That is the whole difficulty.

\section{Why the fitter missed it}

\subsection{The branch that was supposed to catch this}

\fn{sectors.series\_descriptor} contains a branch, \fn{\_volume\_fraction\_base},
whose entire job is this case: it asks whether the residual $D_{\max}/2^N$ is
better explained by a power of $N$ (in which case the base is exactly $2$ and
$\alpha$ comes from that fit) or by an exponential in $N$ (in which case the base
is genuinely below $2$). It requires the power model to win on the last six
points \emph{and} on the full series, and refuses if the series falls back down
or if $|\alpha|>8$.

It refused here, and it refused by a hair. Table~\ref{tab:below} reports the
relative margin of the tail comparison:

\begin{center}
\small
\iftable{tab_below_curve_obc0.tex}
\label{tab:below}
\end{center}

\noindent Group A loses the six-point tail test by $1.3\%$ in residual sum ---
$0.00682$ against $0.00673$ --- while \emph{winning} the full-series test by
$18.9\%$. $W_{74}$ loses the tail by $0.2\%$ and $W_{229}$ by $1.0\%$. Only
$W_{88}$ loses both, and it loses the full test by $14.7\%$. Five of the seven
would have been caught by a rule that asked the full series alone.

The reason the two windows disagree is a genuine near-degeneracy, not a bug:
\begin{equation}
c\cdot 2^N N^{-0.41} \quad\text{and}\quad 1.9226^N
\end{equation}
agree to better than $2\%$ over $N=6\ldots16$, because $\ln N$ is very nearly
affine in $N$ over a window that short. Loosening the tail test until it accepted
group A would be tuning a threshold to a known answer, and it would then have to
be defended against every other rule in the sweep.

\subsection{The certified control settles that no estimator would do better}

The decisive check is a rule whose base is known by derivation rather than by
fit. $W_{134}$'s $D_{\max}$ is exactly rule $150$'s central binomials
$924, 1716, 3003, 6435, 12870, 24310,\ldots$, i.e.\ $\binom{N}{\lfloor N/2\rfloor}
\sim 2^N/\sqrt{\pi N/2}$, so its base is $2$ with $\alpha=-1/2$ and there is
nothing to discover. Fitted with a pure exponential over the same window it
returns $1.9244$, and its rolling six-point base runs $1.904$--$1.938$ ---
\emph{inside} the range $1.858$--$2.039$ that the seven span over the same
window. There is no residual signal in the fit that separates the two cases.

\begin{center}
\small
\iftable{tab_below_controls_obc0.tex}
\end{center}

\noindent The complementary control is $W_{28}$, whose base really is below $2$:
it sits flat at $1.68$--$1.73$ and never approaches the band. So the estimator is
not broken --- it separates a real sub-$2$ base from base $2$ perfectly well ---
it simply cannot separate $2$-with-a-power-law-prefactor from $1.92$ on $N\le16$.
Deriving $b$ is not a workaround for a weak fit; it is the only thing that can
settle the question at these sizes.

\section{The extended sweep to $N=22$}
\label{sec:extended}

The seven were re-run to $N=22$ (background sweep, \fn{wcc\_sweep --lax}), giving
a lever arm two-thirds longer than the uniform window. Three things came out of
it, and one of them cuts against the tidiest version of the story.

\paragraph{The staircases survive.} Both closed forms hold at every new $N$. This
is the load-bearing fact --- everything else rests on $a=1$ --- and it is now
checked over $17$ system sizes rather than $11$.

\paragraph{Group A's fitted base climbs like the control's.} The pure-exponential
base fitted on $N=6\ldots M$ rises monotonically with $M$ for $W_6$, $W_{14}$,
$W_{20}$, $W_{84}$: $1.8951$ at $M=10$ to $1.9345$ at $M=22$, tracking the
analytically continued central binomial ($1.8867\to1.9291$) almost step for step,
while a synthetic series with a genuine base $1.9226$ is flat at $1.9226$ by
construction. That is the signature of a power-law prefactor on base $2$.

\paragraph{Group B's does not, and we should not pretend otherwise.} For
$W_{74}$, $W_{88}$ and $W_{229}$ the same estimator wanders or drifts
\emph{down}: $W_{74}$ goes $1.9884\to1.9313$. The cause is visible in the data
--- their $n_\wcc$ staircase has period $3$ and $D_{\max}$ inherits a strong
modulation from it ($W_{74}$'s residual alternates by nearly $50\%$ between odd
and even $N$) --- and a base fitted across residue classes absorbs part of the
modulation. The climbing argument is therefore corroboration for group A only.
Panel (c) of Fig.~\ref{fig:below} is titled accordingly.

\paragraph{A model-free refutation, for one rule.} Restricting to a residue class
removes the modulation and enables a diagnostic that needs no model at all. For a
pure exponential $c\,b^N$, \emph{every} secant
\[
\left(\frac{D_{\max}(N_2)}{D_{\max}(N_1)}\right)^{1/(N_2-N_1)},
\qquad N_1\equiv N_2 \pmod m ,
\]
equals $b$ exactly. All seven exceed their own fitted base over the tail. And
$W_{229}$'s odd branch \emph{stops decaying} relative to $2^N$ at $N=11$ and then
rises --- $D_{\max}/2^N = 0.3867, 0.3965, 0.4219, 0.4341, 0.4355, 0.4286$ at
$N=11,13,15,17,19,21$ --- giving a secant of $\mathbf{2.0052}$ over
$N=15\ldots21$. A base below $2$ forces $D_{\max}/2^N=c(b/2)^N$ to decay
geometrically forever, so this single branch refutes the fitted model for
$W_{229}$ from the data alone, with no asymptotics and no theorem. The same
diagnostic returns $1.7321$ for $W_{28}$, so it is not firing on noise.

\begin{figure}[htbp]
\centering
\ifgraphic{../../figures/fig_below_curve_obc0.pdf}{\textwidth}
\caption{\textbf{Two ways to sit under the curve.} (a) The corner of R9's F1.
The green cell is $40$ V-free rules at exactly $(1,2)$ --- on the curve --- drawn
at $a=0.981$ where the curve is $2.0387$, and so \emph{below} it: a rendering
artefact (\S\ref{sec:artefact}). The open red markers are the seven as fitted;
the arrow is the correction to the derived $(1,2)$. (b) The residual
$D_{\max}/2^N$ on log--log axes, normalised at $N=6$, even branch: a power-law
prefactor is a straight line, a genuine base bends. The certified control
$W_{134}$ (blue squares, base $2$ by derivation) and the null (grey dashed, a
genuine base $1.9226$) bracket the question, and over this window they barely
separate. The purple branch is $W_{229}$ at \emph{odd} $N$, which stops decaying
--- the model-free refutation of \S\ref{sec:extended}. (c) The pure-exponential
base fitted on $N=6\ldots M$ against $M$. The null is flat by construction; the
continued central binomial climbs; group A climbs with it; group B wanders,
because its period-$3$ modulation contaminates a base fitted across residues.}
\label{fig:below}
\end{figure}

\section{The correction}

Seven \fn{ANALYTIC} entries were installed in
\fn{scaling/sectors.py}, keyed \verb|(rule, "obc0")|:
\[
n_\wcc:\;(a,\alpha)=(1,1),\qquad D_{\max}:\;b=2,\;\alpha=\texttt{None}.
\]
The \verb|None| is deliberate and required a small extension to
\fn{series\_descriptor}. An analytic entry may derive the \emph{base} without
deriving the prefactor power that goes with it; writing a fitted number into a
table headed ``derived rather than fitted'' would be worse than refitting it
openly. When $\alpha$ is \verb|None| the power is now refitted with the derived
base held fixed, via the existing \fn{\_alpha\_at\_base}, and the descriptor
records \fn{alpha\_source = "refit at the derived base"}. The results are
$\alpha=-0.413$ (group A), $-0.295$ ($W_{74}$), $-0.326$ ($W_{88}$), $-0.437$
($W_{229}$) --- all comfortably inside the $|\alpha|\le8$ guard and all of the
size a prefactor power should be.

\paragraph{Effect on the sector map.}
\begin{center}
\begin{tabular}{lrr}
\toprule
verdict & before & after \\
\midrule
on the curve, $ab=2$ & $197$ & $204$ \\
above, $ab>2$ & $50$ & $50$ \\
sub-exponential factor, test degenerate & $7$ & $0$ \\
irregular series, no growth law & $2$ & $2$ \\
\midrule
\textbf{rules with $ab<2$} & $\mathbf{7}$ & $\mathbf{0}$ \\
\bottomrule
\end{tabular}
\end{center}

\noindent The anchors are untouched: $W_{204}$, $W_{51}$ and $W_{150}$ remain at
$ab=2.000000$ and the six room-packing rules remain at
$\varphi\cdot4^{1/5}=2.1350$. The full validation pass reports
\verb|hyperbola_violations: 0| and \verb|sum_rule_max_abs_residual: 0| over
$2926$ rule--$N$ units.

\paragraph{Corroborating structure.} The hyperbola is tight when the sectors are
comparable in size, which is what these rules look like. $W_{14}$ at $N=16$ has nine
sectors of sizes $23234, 18250, 13456, 5666, 3830, 550, 511, 38, 1$, summing to
$2^{16}$: the top five are within a factor of six of each other and carry $98.3\%$
of the basis. That is the equality condition of \eqref{eq:hyper} --- comparable
sectors, a bounded number of them --- not the profile of a rule with a smaller
base.

\section{Answering the question as asked}

\begin{enumerate}
\item \textbf{Which rules?} Seven have a genuinely sub-$2$ fitted product:
$W_6$, $W_{14}$, $W_{20}$, $W_{74}$, $W_{84}$, $W_{88}$, $W_{229}$, in two groups
of four and three sharing identical series. Separately, $40$ V-free rules that
are exactly on the curve are \emph{drawn} below it by F1's family offset; that
marker is larger than the real effect and is the first thing a reader sees.

\item \textbf{Is it simply a fitting problem?} Yes --- in the precise sense that
no sector count and no sector size is wrong, the finite-$N$ bound
\eqref{eq:finite} holds at every $N$ for all $256$ rules, and only the summary
pair $(a,b)$ was wrong. But ``simply'' would be misleading if it suggested a
better fit is the remedy. The certified control shows that on $N\le16$ (indeed on
$N\le22$) no pure-exponential estimator can separate base $2$ with an $N^{-0.4}$
prefactor from a genuine base $1.92$. The remedy is the theorem: $a=1$ is exact
from the staircase, and $a=1$ forces $b=2$.

\item \textbf{What changed?} Seven derived overrides, one new refit path for
analytic entries that fix a base but not a power, a corrected note in F1's
caption and header, and a sector map with nothing below the curve.
\end{enumerate}

\section{Reproduction}

\begin{verbatim}
# the analysis, its figure and its two tables
python -m qca_fragmentation.scaling.below_curve --bc obc0 --rebuild

# the sector map and R9's figures/tables, which consume the overrides
python -m qca_fragmentation.scaling.sector_figure --bc obc0 --rebuild
python -m qca_fragmentation.scaling.r9_tables     --bc obc0

# the extension sweep behind Section 5 (already on disk to N=22)
python -m qca_fragmentation.wcc_sweep --rules 6,14,20,74,84,88,229 \
       --bc obc0 --n-min 17 --n-max 22 --lax

# the regression tests
python -m pytest qca_fragmentation/tests/test_below_curve.py -q
\end{verbatim}

\noindent \fn{scaling/below\_curve.py} is written so that it can still explain
what it corrected after the correction is in place: \fn{sectors.R16\_OVERRIDES}
names exactly the entries R16 added, and the context manager
\fn{without\_r16\_override} lifts them so the pre-correction fit is recomputed
from the series rather than hard-coded into the prose. Every number in this
report comes out of \fn{analytics/below\_curve\_obc0.json}.

The regression suite pins the parts that could rot silently: that
\fn{R16\_OVERRIDES} still matches the set of rules actually below the curve, that
lifting the override is reversible, that both closed forms still hold at every
$N$, that \fn{below\_curve.FAM\_DX} still equals \fn{sector\_figure.\_FAM\_DX}
(otherwise \S\ref{sec:artefact} would be explaining a figure that no longer
exists), that group A still climbs and group B still does not, and that the
$\alpha$ overrides parse under matplotlib's mathtext as well as under \LaTeX{}
--- a test which caught a real defect, \verb|\ge| being valid \LaTeX{} but not
valid mathtext.

\end{document}
